\documentclass{beamer}

\usepackage{pifont}
\usepackage{tikz}
\usetikzlibrary{arrows.meta,positioning, calc}
\usepackage{caption}
\usepackage[ngerman]{babel}
\usepackage{csquotes}

\newcommand{\checkicon}{\ding{51}}
\newcommand{\warnicon}{\ding{115}}

\usetheme{metropolis}
\title{Offene Daten erlebbar machen}
\date{}
\author{Adaktylos Leonidas Odysseas, Baumgartner Vincent Vitez,\\ Neda Nikkhousarighiyeh, Alexander Grath}
\begin{document}

  % ============================================================

  % ============================================================
  \begin{frame}{Erster Eindruck der Testpersonen}
    \begin{itemize}
      \setlength\itemsep{0.8em}
      \item App wurde sofort mit Google Maps verglichen
      \item Konzept kam gut an – Nutzen war schnell klar
      \item Technisch versierte TPs flogen durch alle Aufgaben
      \item Swipe-Geste war die einzige universelle Hürde
    \end{itemize}
  \end{frame}

  % ============================================================
   \begin{frame}{SUS-Scores}
    \centering
    \begin{tabular}{@{}lll@{}}
      \textbf{TP} & \textbf{Score} & \textbf{Bewertung} \\
      \hline
      TP1 & 72{,}5 & gut \\
      TP2 & 62{,}5 & akzeptabel \\
      TP3 & 80{,}0 & exzellent \\
      TP4 & 77{,}5 & gut \\
      TP5 & 60{,}0 & akzeptabel \\
      TP6 & 70{,}0 & gut \\
      TP7 & 87{,}5 & exzellent \\
      TP8 & 85{,}0 & exzellent \\
      \hline
      \textbf{Ø} & \textbf{74{,}4} & \textbf{gut} \\
    \end{tabular}

    \vspace{0.4cm}
    \small Technisch versierte TPs hatten kaum Probleme. App wurde durchgehend mit Google Maps verglichen.
  \end{frame}

  % ============================================================
  \begin{frame}{Erkenntnisse}
    \textbf{Stärken}
    \begin{itemize}
      \item Farbkodierte Pins intuitiv
      \item Suche funktioniert für alle
      \item Konzept als nützlich bewertet
    \end{itemize}

    \vspace{0.4cm}
    \textbf{Schwächen}
    \begin{itemize}
      \item Swipe-to-reveal nicht entdeckbar
      \item Pin-Highlight nach Suche fehlt
      \item Download-Bestätigung fehlt
    \end{itemize}
  \end{frame}

  % ============================================================
  \begin{frame}{Reflection}
    \begin{itemize}
      \setlength\itemsep{0.7em}
      \item Faire Aufgabenteilung über alle Meilensteine
      \item GitHub + Discord haben Kollaboration einfach gehalten
      \item \textbf{Frühzeitiges Nutzerfeedback} hätte viele Probleme früher aufgedeckt
      \item \textbf{UX ist mehr als Technik} – Nutzerfreundlichkeit braucht Iteration
    \end{itemize}
  \end{frame}

  % ============================================================
  \begin{frame}{Schlussfolgerung \& Ausblick}
    Dat-O-Kart füllt eine reale Lücke bei mobilen Open-Data-Apps für Wien.

    \vspace{0.5cm}
    \textbf{Nächste Schritte}
    \begin{itemize}
      \item Live-Anbindung an \textit{data.wien.gv.at}
      \item Filteroptionen für Pins
      \item Swipe durch sichtbares UI-Element ersetzen
    \end{itemize}
  \end{frame}


 \end{document}
